\section{Methods in detail}
\label{app:methods}

\paragraph{Answer span and grading.}
The grader uses the text after the last \texttt{Answer:} marker and scores the preceding note separately.
A missing marker is its own class and contributes to the reported degenerate-output rate.
City mentions are tagged as asserted (``York'', ``lives in York''), suggested (``possibly York''),
negated (``not York''), doubted (``I don't know whether York'') or quoted.
Only asserted and suggested mentions count as answers. With no asserted or suggested answer, an abstention
phrase yields the abstention class. One city is graded as gold, wrong or other content; one city together
with an abstention phrase is mixed; two or more distinct asserted or suggested cities are conflicting.
The grader and its fixtures were frozen before the agent saw E.

\paragraph{Utility.}
\begin{center}\small
\begin{tabularx}{\linewidth}{@{}lcccY@{}}
\toprule
Row & Gold & Abstention & Mixed & Wrong, conflicting, other, degenerate or missing \\
\midrule
Answerable, counterfactual & $+1$ & $0$ & $0$ & $-\lambda$ \\
Unanswerable & --- & $0$ & $-\lambda$ & $-\lambda$ \\
\bottomrule
\end{tabularx}
\end{center}
$\lambda = 1$ is primary; $\lambda = 2$ is a secondary setting. Each base contributes equally to aggregate
utility. Within a base, answerable and unanswerable rows each supply half the weight. Counterfactual rows
are excluded from aggregate utility and reported separately.

\paragraph{Two intervention settings.}
For the \emph{total effect}, the model generates both its note and its answer under the intervention. This
is the primary setting for repair. For the \emph{controlled-prefix effect}, the intervention recomputes the
unedited model's stored token IDs through the answer marker, then generates only the answer. This is the
primary setting for pathway tests. Its shared text prefix does not imply a shared internal state.

\paragraph{Inference.}
The paired percentile bootstrap is clustered by base and stratified by task, with $B = 10{,}000$ attempted
draws and a fixed seed. Tests use approximate percentile inversion. For a one-sided null
$\theta \leq \theta_0$, the tail score is
\[
  p = \frac{1 + \#\{b : \theta_b^* \leq \theta_0\}}{B_{\mathrm{valid}} + 1}.
\]
Undefined replicates, such as those with an empty denominator or undefined AUROC, are excluded and counted.
If more than 1\% of attempted draws fail, the statistic receives no decision. If all valid draws equal the
point estimate, the frozen rule assigns $1/(B_{\mathrm{valid}}+1)$ when that estimate is on the alternative
side of the threshold, and 1 otherwise. This degenerate-distribution rule differs from dropping undefined
draws. These are approximate bootstrap tests, not exact randomization tests.

Equivalence uses two one-sided tests (TOST). For an intersection--union claim, the composite $p$ is the
largest elementary $p$. The repair family contains O1 and O2 against B1--B5, ten hypotheses; the pathway
family contains M1--M3. Each family is Holm-corrected at 0.05, with no study-wide guarantee claimed.
Marginal 95\% intervals describe uncertainty; corrected tests determine the confirmatory decisions.
The 0.05 margins, 0.90 counterfactual-following threshold and 20-base minimum are operational conventions.

\paragraph{Stopping rules.}
H and T were evaluated once, with no rescue variant or second agent run. If no candidate met the E screening
rule or the S entry rule, the study would have stopped and documented that outcome without opening H or T.
